\section{Regolazione della ricerca predittiva}
\label{sec:regulator}
\subsection{Quando cambiare candidato}
Il metodo C usa le stesse previsioni e potature di B. Il regolatore osserva
se la ricerca continua senza migliorare il giro e, in quel caso, prova a
cambiare il prossimo candidato da approfondire. Un cambiamento di questo
tipo viene chiamato \emph{diversificazione}.

Per decidere quando intervenire uso due contatori:
\begin{itemize}
\item $q$: quanti stati sono stati estratti dall'ultimo miglioramento
  della distanza del giro migliore;
\item $c$: quante estrazioni ordinarie sono avvenute dall'ultima
  diversificazione, per lasciare tempo alla ricerca predittiva fra due interventi.
\end{itemize}
Entrambi partono da zero. Fisso una soglia, indicata con $P$, uguale a
$100$. Il regolatore interviene soltanto quando \emph{entrambi} i contatori
hanno raggiunto la soglia:
\[
 q\geq P\quad\text{e}\quad c\geq P.
\]
Quando la condizione è soddisfatta, invece di prendere l'ultimo candidato
inserito nella pila scelgo quello con la distanza già percorsa $g$ minore.
In caso di parità scelgo il primo presente nella pila. Le altre alternative
restano disponibili. Se la condizione non è soddisfatta, la scelta rimane
quella ordinaria di B.

La corrispondenza con la consegna è quindi: le \emph{osservazioni} sono i
due contatori; la \emph{regola} verifica le soglie; il \emph{controllo}
cambia quale candidato riceve il prossimo passo di ricerca.

\subsection{Come cambiano i contatori}
Ogni estrazione aumenta $q$, anche se il candidato è un giro completo o
viene subito scartato. Trovare un giro strettamente più corto del migliore
lo azzera. Prima del primo giro, $q$ conta le estrazioni dall'avvio.

Un'estrazione ordinaria aumenta $c$. Una diversificazione lo azzera,
così devono trascorrere almeno altre cento estrazioni ordinarie prima di
un nuovo intervento. La condizione viene controllata prima di estrarre:
se dall'avvio non ci sono miglioramenti, il primo intervento può quindi
avvenire al passo $101$.

\subsection{Motivazione e limiti}
L'assenza prolungata di miglioramenti è chiamata \emph{stagnazione}.
È un segnale semplice: suggerisce di provare un candidato diverso,
ma non dimostra che il modello stia sbagliando. Anche scegliere il costo
parziale minore può essere poco utile, perché un percorso appena iniziato
ha spesso una distanza bassa pur lasciando molto lavoro da fare.
Il nuovo candidato può inoltre appartenere allo stesso primo ramo.

Per trovare il minimo devo esaminare tutti i candidati nella pila.
Se sono molti, questo intervento richiede tempo; espandere percorsi poco
profondi può anche aggiungere molti nuovi candidati. Le prove misurano
quindi sia la qualità dei giri sia il lavoro svolto.

La regola e la soglia restano fisse durante le prove. Non ho cercato
sistematicamente la soglia migliore sulle mappe finali e non alleno il
regolatore durante la ricerca. Il confronto B--C serve proprio a verificare
quando questa scelta semplice aiuta e quando peggiora il risultato.
